\ifdefined\UOMTransportFragmentLoaded\else
\documentclass[11pt]{article}
\usepackage[margin=1in]{geometry}
\usepackage{amsmath,amssymb,amsthm,mathtools}
\usepackage{bm}
\usepackage{microtype}
\usepackage{enumitem}
\usepackage{booktabs}
\usepackage{array}
\usepackage{xcolor}
\usepackage{xr-hyper}
\usepackage{hyperref}
\usepackage[nameinlink,noabbrev]{cleveref}

\providecommand{\Fsloc}{\mathcal F_{\rm s.loc}}
\providecommand{\Qt}{Q_t}
\externaldocument[main-]{UOM/main}
\externaldocument[uomdesc-]{UOM/uom_descendant_note}
\externaldocument{UOM/completion_sector_split_model_note}

\hypersetup{
  colorlinks=true,
  linkcolor=blue!60!black,
  citecolor=blue!60!black,
  urlcolor=blue!60!black
}

\theoremstyle{definition}
\newtheorem{definition}{Definition}[section]
\newtheorem{assumption}{Assumption}[section]
\theoremstyle{plain}
\newtheorem{theorem}[definition]{Theorem}
\newtheorem{lemma}[definition]{Lemma}
\newtheorem{proposition}[definition]{Proposition}
\newtheorem{corollary}[definition]{Corollary}
\theoremstyle{remark}
\newtheorem{remark}[definition]{Remark}
\crefname{definition}{definition}{definitions}
\Crefname{definition}{Definition}{Definitions}
\crefname{assumption}{assumption}{assumptions}
\Crefname{assumption}{Assumption}{Assumptions}
\crefname{theorem}{theorem}{theorems}
\Crefname{theorem}{Theorem}{Theorems}
\crefname{lemma}{lemma}{lemmas}
\Crefname{lemma}{Lemma}{Lemmas}
\crefname{proposition}{proposition}{propositions}
\Crefname{proposition}{Proposition}{Propositions}
\crefname{corollary}{corollary}{corollaries}
\Crefname{corollary}{Corollary}{Corollaries}
\crefname{remark}{remark}{remarks}
\Crefname{remark}{Remark}{Remarks}

\newcommand{\cH}{\mathcal H}
\newcommand{\cA}{\mathcal A}
\newcommand{\cD}{\mathcal D}
\newcommand{\cB}{\mathcal B}
\newcommand{\cC}{\mathcal C}
\newcommand{\cL}{\mathcal L}
\newcommand{\cK}{\mathcal K}
\newcommand{\cR}{\mathcal R}
\newcommand{\cG}{\mathcal G}
\newcommand{\Ztwo}{\mathbb Z_2}
\newcommand{\Sys}{\mathrm{Sys}}
\newcommand{\Time}{\mathrm{Time}}
\newcommand{\RG}{\mathrm{RG}}
\newcommand{\Tr}{\mathrm{Tr}}
\newcommand{\diag}{\mathrm{diag}}
\newcommand{\ad}{\mathrm{ad}}
\newcommand{\id}{\mathrm{id}}
\newcommand{\norm}[1]{\left\lVert #1\right\rVert}
\newcommand{\abs}[1]{\left\lvert #1\right\rvert}
\newcommand{\ip}[2]{\left\langle #1,#2\right\rangle}

\externaldocument{UOM/compressed_log_spectral_selection_note}
\externaldocument{UOM/canonical_odd_interface_note}
\externaldocument{UOM/active_response_transport_program_note}
\title{Raw Odd Selector and Scheme-Stable Affine Output}
\author{}
\date{}
\begin{document}
\maketitle
\fi

\section{Raw Odd Selector and Scheme-Stable Affine Output}
\label{sec:odd-selector}

This section isolates the raw selector step on the odd descendant side of the UOM-to-Main
interface used in Definition~\ref{main-def:uom-main-theorem-assumptions}.
At fixed finite cutoff, it records the \(\Gamma\)-odd affine output determined by a chosen
sufficiently local anomaly representative and its same-scheme and inter-scheme translation laws.

\paragraph{Convention on the two \texorpdfstring{\(\mathbb Z_2\)}{Z2} gradings.}
The cohomological degree for the \(Q_t\)-complex is written with bars:
\[
(\cdot)_{\bar 0},
\qquad
(\cdot)_{\bar 1}.
\]
Physical parity refers to the swap involution \(\Gamma\), with eigenspaces
\[
\Gamma F=\pm F.
\]
Thus the selector output constructed below is cohomologically degree \(0\) but physically odd.
Unless bar-degree notation is used explicitly, ``odd'' below means \(\Gamma\)-odd rather than
\(Q_t\)-degree \(1\).

Work at one accepted tick \(k\) and one branch value \(u\) at cutoff \(\Lambda_k\).
Use the raw UOM descendant package of
\cref{uomdesc-thm:uomdesc-raw-package,uomdesc-thm:uomdesc-conditional-odd-package}.
In particular:
\begin{enumerate}[leftmargin=2em]
\item a sufficiently local good monotone representative
  \[
  \mathcal O^{(0)}_{\Lambda_k}(u);
  \]
\item its canonical anomaly representative
  \[
  \mathcal A^{(1)}_{\Lambda_k}(u)
  :=
  Q^{\rm phys}_{u,\Lambda_k}\mathcal O^{(0)}_{\Lambda_k}(u);
  \]
\item a swap involution
  \[
  \Gamma:
  \big(\mathcal F_{\rm s.loc}^{\Lambda_k}(\mathcal X_k)\big)_{\bar\imath}
  \to
  \big(\mathcal F_{\rm s.loc}^{\Lambda_k}(\mathcal X_k)\big)_{\bar\imath}
  \qquad
  (\bar\imath\in\{\bar 0,\bar 1\})
  \]
  preserving sufficient locality and cohomological degree, and commuting with \(Q_t\):
  \[
  \Gamma Q_t=Q_t\Gamma;
  \]
\item swap-evenness of the canonical core,
  \[
  \Gamma\mathcal O^{(0)}_{\Lambda_k}(u)=\mathcal O^{(0)}_{\Lambda_k}(u),
  \qquad
  \Gamma\mathcal A^{(1)}_{\Lambda_k}(u)=\mathcal A^{(1)}_{\Lambda_k}(u);
  \]
\item swap covariance of the fixed-scheme physical radial differential on the chosen monotone
  sector,
  \[
  \Gamma Q^{\rm phys}_{u,\Lambda_k}
  =
  Q^{\rm phys}_{u,\Lambda_k}\Gamma;
  \]
\end{enumerate}
Items (1) and (2) are contained in \cref{uomdesc-thm:uomdesc-raw-package};
item (3) uses \cref{uomdesc-def:uomdesc-gamma,uomdesc-lem:uomdesc-gamma-comm};
item (4) uses \cref{uomdesc-prop:uomdesc-gamma-even-core}; and item (5) uses
\cref{uomdesc-prop:uomdesc-swap-split}.

Fix a chosen \(\Gamma\)-stable selector class of sufficiently local odd representatives of the same
\(Q_t\)-cohomology class as \(\mathcal A^{(1)}_{\Lambda_k}(u)\), and restrict to representatives
for which the exact descent equation
\begin{equation}
Q^{\rm phys}_{u,\Lambda_k}\mathcal O^{(0)}_{\Lambda_k}(u)
+
Q_t\widetilde{\Xi}^{(1)}_{\Lambda_k}(u)
=
\widetilde{\mathcal A}^{(1)}_{\Lambda_k}(u)
\label{eq:selector-full-descent}
\end{equation}
admits at least one sufficiently local even solution
\(
\widetilde{\Xi}^{(1)}_{\Lambda_k}(u)
\in
\big(\mathcal F_{\rm s.loc}^{\Lambda_k}(\mathcal X_k)\big)_{\bar 0}
\).

\begin{definition}[Selector class and canonical odd kernel channel]
\label{def:selector-class}
Let
\[
\mathcal C_{{\rm sel},\Lambda_k}(u)
\subset
\big(\mathcal F_{\rm s.loc}^{\Lambda_k}(\mathcal X_k)\big)_{\bar 1}
\]
be the chosen \(\Gamma\)-stable class of sufficiently local odd representatives of the same
degree-one \(Q_t\)-cohomology class as \(\mathcal A^{(1)}_{\Lambda_k}(u)\) for which
\eqref{eq:selector-full-descent} admits a sufficiently local even solution.
For
\(
\widetilde{\mathcal A}^{(1)}_{\Lambda_k}(u)\in\mathcal C_{{\rm sel},\Lambda_k}(u)
\),
define its swap-even and swap-odd parts by
\[
\widetilde{\mathcal A}^{(1)}_{\Lambda_k,\pm}(u)
:=
\frac12\Big(
\widetilde{\mathcal A}^{(1)}_{\Lambda_k}(u)
\pm
\Gamma\widetilde{\mathcal A}^{(1)}_{\Lambda_k}(u)
\Big).
\]
Define the canonical odd kernel channel by
\[
\mathcal K^-_{\Lambda_k}
:=
\ker(Q_t)
\cap
\big(\mathcal F_{\rm s.loc}^{\Lambda_k}(\mathcal X_k)\big)_{\bar 0}
\cap
\ker(\Gamma+\mathrm{id}).
\]
\end{definition}

\begin{proposition}[Odd defect of a selector representative is always \(Q_t\)-exact]
\label{prop:selector-odd-defect}
Let
\(
\widetilde{\mathcal A}^{(1)}_{\Lambda_k}(u)\in\mathcal C_{{\rm sel},\Lambda_k}(u)
\).
Then:
\begin{enumerate}[leftmargin=2em]
\item
  \(
  \widetilde{\mathcal A}^{(1)}_{\Lambda_k,+}(u)
  \)
  and
  \(
  \widetilde{\mathcal A}^{(1)}_{\Lambda_k,-}(u)
  \)
  are sufficiently local;
\item
  \(
  \widetilde{\mathcal A}^{(1)}_{\Lambda_k,+}(u)
  \)
  is a swap-even representative of the same \(Q_t\)-cohomology class as
  \(\mathcal A^{(1)}_{\Lambda_k}(u)\);
\item
  \[
  \widetilde{\mathcal A}^{(1)}_{\Lambda_k,-}(u)
  \in
  \operatorname{im}(Q_t)
  \cap
  \big(\mathcal F_{\rm s.loc}^{\Lambda_k}(\mathcal X_k)\big)_{\bar 1}
  \cap
  \ker(\Gamma+\mathrm{id}).
  \]
\end{enumerate}
\end{proposition}

\begin{proof}
Sufficient locality is preserved by linear combinations, proving item (1).

Because
\(
\widetilde{\mathcal A}^{(1)}_{\Lambda_k}(u)
\)
and
\(
\mathcal A^{(1)}_{\Lambda_k}(u)
\)
represent the same \(Q_t\)-cohomology class, there exists
\(
Y_{\Lambda_k}(u)\in
\big(\mathcal F_{\rm s.loc}^{\Lambda_k}(\mathcal X_k)\big)_{\bar 0}
\)
such that
\[
\widetilde{\mathcal A}^{(1)}_{\Lambda_k}(u)
-
\mathcal A^{(1)}_{\Lambda_k}(u)
=
Q_tY_{\Lambda_k}(u).
\]
Taking the swap-even part and using
\(
\Gamma Q_t=Q_t\Gamma
\)
and
\(
\Gamma\mathcal A^{(1)}_{\Lambda_k}(u)=\mathcal A^{(1)}_{\Lambda_k}(u)
\),
we obtain
\[
\widetilde{\mathcal A}^{(1)}_{\Lambda_k,+}(u)
-
\mathcal A^{(1)}_{\Lambda_k}(u)
=
Q_tY_{\Lambda_k,+}(u),
\]
so item (2) follows.

Taking the swap-odd part of the same identity gives
\[
\widetilde{\mathcal A}^{(1)}_{\Lambda_k,-}(u)
=
Q_tY_{\Lambda_k,-}(u).
\]
By construction,
\(
\Gamma\widetilde{\mathcal A}^{(1)}_{\Lambda_k,-}(u)
=
-\widetilde{\mathcal A}^{(1)}_{\Lambda_k,-}(u)
\),
which proves item (3).
\end{proof}

\begin{definition}[Selected odd descendant affine class]
\label{def:selector-affine-class}
For
\(
\widetilde{\mathcal A}^{(1)}_{\Lambda_k}(u)\in\mathcal C_{{\rm sel},\Lambda_k}(u)
\),
define the associated raw odd selector output by
\[
\mathfrak X_{\Lambda_k}^{{\rm sel},-}
\big(u;\widetilde{\mathcal A}^{(1)}_{\Lambda_k}\big)
:=
\left\{
\Xi^-_{\Lambda_k}(u)
\in
\big(\mathcal F_{\rm s.loc}^{\Lambda_k}(\mathcal X_k)\big)_{\bar 0}
\;\middle|\;
\Gamma\Xi^-_{\Lambda_k}(u)=-\Xi^-_{\Lambda_k}(u),
\quad
Q_t\Xi^-_{\Lambda_k}(u)
=
\widetilde{\mathcal A}^{(1)}_{\Lambda_k,-}(u)
\right\}.
\]
\end{definition}

\begin{remark}[Dependence only on the swap-odd defect]
\label{rem:selector-odd-defect-only}
By \cref{thm:selector-decomposition}, the affine class
\(
\mathfrak X_{\Lambda_k}^{{\rm sel},-}
\big(u;\widetilde{\mathcal A}^{(1)}_{\Lambda_k}\big)
\)
depends only on the swap-odd defect
\(
\widetilde{\mathcal A}^{(1)}_{\Lambda_k,-}(u)
\).
\end{remark}

\begin{theorem}[Selector decomposition into even representative channel and odd output]
\label{thm:selector-decomposition}
Let
\(
\widetilde{\mathcal A}^{(1)}_{\Lambda_k}(u)\in\mathcal C_{{\rm sel},\Lambda_k}(u)
\),
and let
\(
\widetilde{\Xi}^{(1)}_{\Lambda_k}(u)
\)
be any sufficiently local even solution of \eqref{eq:selector-full-descent}.
Define
\[
\widetilde{\Xi}^{(1)}_{\Lambda_k,\pm}(u)
:=
\frac12\Big(
\widetilde{\Xi}^{(1)}_{\Lambda_k}(u)
\pm
\Gamma\widetilde{\Xi}^{(1)}_{\Lambda_k}(u)
\Big).
\]
Then:
\begin{enumerate}[leftmargin=2em]
\item the swap-even and swap-odd channels decouple as
  \begin{align}
  Q^{\rm phys}_{u,\Lambda_k}\mathcal O^{(0)}_{\Lambda_k}(u)
  +
  Q_t\widetilde{\Xi}^{(1)}_{\Lambda_k,+}(u)
  &=
  \widetilde{\mathcal A}^{(1)}_{\Lambda_k,+}(u),
  \label{eq:selector-even-channel}
  \\
  Q_t\widetilde{\Xi}^{(1)}_{\Lambda_k,-}(u)
  &=
  \widetilde{\mathcal A}^{(1)}_{\Lambda_k,-}(u);
  \label{eq:selector-odd-channel}
  \end{align}
\item the raw odd selector output
  \(
  \mathfrak X_{\Lambda_k}^{{\rm sel},-}
  \big(u;\widetilde{\mathcal A}^{(1)}_{\Lambda_k}\big)
  \)
  contains the odd projection of any full solution; in particular it is nonempty;
\item it is an affine space over the canonical odd kernel channel:
  \[
  \mathfrak X_{\Lambda_k}^{{\rm sel},-}
  \big(u;\widetilde{\mathcal A}^{(1)}_{\Lambda_k}\big)
  =
  \widetilde{\Xi}^{(1)}_{\Lambda_k,-}(u)
  +
  \mathcal K^-_{\Lambda_k}.
  \]
\end{enumerate}
\end{theorem}

\begin{proof}
Apply \(\Gamma\) to \eqref{eq:selector-full-descent}.
Using
\(
\Gamma Q^{\rm phys}_{u,\Lambda_k}
=
Q^{\rm phys}_{u,\Lambda_k}\Gamma
\)
on the chosen monotone sector and
\(
\Gamma\mathcal O^{(0)}_{\Lambda_k}(u)=\mathcal O^{(0)}_{\Lambda_k}(u)
\),
we obtain
\[
Q^{\rm phys}_{u,\Lambda_k}\mathcal O^{(0)}_{\Lambda_k}(u)
+
Q_t\Gamma\widetilde{\Xi}^{(1)}_{\Lambda_k}(u)
=
\Gamma\widetilde{\mathcal A}^{(1)}_{\Lambda_k}(u).
\]
Adding and subtracting this equation from \eqref{eq:selector-full-descent} yields
\eqref{eq:selector-even-channel} and \eqref{eq:selector-odd-channel}, proving item (1).

Item (2) is immediate from item (1): the odd projection
\(
\widetilde{\Xi}^{(1)}_{\Lambda_k,-}(u)
\)
of any full solution belongs to the set in
\cref{def:selector-affine-class}.

Now let
\(
\Xi^-_{\Lambda_k,1}(u)
\)
and
\(
\Xi^-_{\Lambda_k,2}(u)
\)
lie in
\(
\mathfrak X_{\Lambda_k}^{{\rm sel},-}
\big(u;\widetilde{\mathcal A}^{(1)}_{\Lambda_k}\big)
\).
Subtracting their defining equations gives
\[
Q_t\big(
\Xi^-_{\Lambda_k,1}(u)-\Xi^-_{\Lambda_k,2}(u)
\big)
=
0.
\]
Because both are \(\Gamma\)-odd, their difference lies in
\(
\mathcal K^-_{\Lambda_k}
\).
Conversely, if
\(
K^-_{\Lambda_k}\in\mathcal K^-_{\Lambda_k}
\),
then
\[
Q_t\big(
\widetilde{\Xi}^{(1)}_{\Lambda_k,-}(u)+K^-_{\Lambda_k}
\big)
=
\widetilde{\mathcal A}^{(1)}_{\Lambda_k,-}(u),
\]
and the sum remains \(\Gamma\)-odd.
This proves item (3).
\end{proof}

\begin{corollary}[The two selector channels]
\label{cor:selector-two-channels}
At fixed cutoff, the selector problem splits into exactly two channels.
\begin{enumerate}[leftmargin=2em]
\item \textbf{Kernel channel.}
  If
  \(
  \widetilde{\mathcal A}^{(1)}_{\Lambda_k,-}(u)=0
  \),
  then
  \[
  \mathfrak X_{\Lambda_k}^{{\rm sel},-}
  \big(u;\widetilde{\mathcal A}^{(1)}_{\Lambda_k}\big)
  =
  \mathcal K^-_{\Lambda_k}.
  \]
  In particular, the canonical anomaly representative
  \(
  \widetilde{\mathcal A}^{(1)}_{\Lambda_k}(u)=\mathcal A^{(1)}_{\Lambda_k}(u)
  \)
  produces only the canonical odd kernel sector.
\item \textbf{Representative channel.}
  If
  \(
  \widetilde{\mathcal A}^{(1)}_{\Lambda_k,-}(u)\neq 0
  \),
  then the raw odd selector output is the translated affine space
  \[
  \widetilde{\Xi}^{(1)}_{\Lambda_k,-}(u)+\mathcal K^-_{\Lambda_k}
  \]
  solving
  \[
  Q_t\Xi^-_{\Lambda_k}(u)
  =
  \widetilde{\mathcal A}^{(1)}_{\Lambda_k,-}(u).
  \]
\end{enumerate}
Thus the selector output is determined by the canonical odd kernel channel together with the
swap-odd part of the chosen anomaly representative.
\end{corollary}

\begin{proof}
This is an immediate specialization of \cref{thm:selector-decomposition}.
\end{proof}

\begin{remark}[Kernel activity versus representative-driven activity]
\label{rem:selector-two-sources}
By \cref{cor:selector-two-channels}, the selector output is either the canonical odd kernel channel
\(\mathcal K^-_{\Lambda_k}\) or a translate of that channel by a particular solution of
\(
Q_t\Xi^-_{\Lambda_k}(u)=\widetilde{\mathcal A}^{(1)}_{\Lambda_k,-}(u)
\).
\end{remark}

\begin{proposition}[Same-scheme affine translation of selector outputs]
\label{prop:selector-same-scheme}
Let
\[
\widetilde{\mathcal A}^{(1)}_{\Lambda_k}(u),
\qquad
\widehat{\mathcal A}^{(1)}_{\Lambda_k}(u)
\in
\mathcal C_{{\rm sel},\Lambda_k}(u)
\]
be selector representatives in the same finite-cutoff scheme, and suppose there exists an even
sufficiently local cochain \(Y_{\Lambda_k}(u)\) such that
\begin{equation}
\widehat{\mathcal A}^{(1)}_{\Lambda_k}(u)
-
\widetilde{\mathcal A}^{(1)}_{\Lambda_k}(u)
=
Q_tY_{\Lambda_k}(u).
\label{eq:selector-same-scheme}
\end{equation}
Define
\[
Y^-_{\Lambda_k}(u)
:=
\frac12\Big(
Y_{\Lambda_k}(u)-\Gamma Y_{\Lambda_k}(u)
\Big).
\]
Then
\[
\mathfrak X_{\Lambda_k}^{{\rm sel},-}
\big(u;\widehat{\mathcal A}^{(1)}_{\Lambda_k}\big)
=
\mathfrak X_{\Lambda_k}^{{\rm sel},-}
\big(u;\widetilde{\mathcal A}^{(1)}_{\Lambda_k}\big)
+
Y^-_{\Lambda_k}(u).
\]
If \(Y'_{\Lambda_k}(u)\) is another primitive satisfying
\eqref{eq:selector-same-scheme}, then
\[
Y_{\Lambda_k}^{\prime -}(u)-Y^-_{\Lambda_k}(u)\in \mathcal K^-_{\Lambda_k}.
\]
In particular, if
\(
\widehat{\mathcal A}^{(1)}_{\Lambda_k,-}(u)
=
\widetilde{\mathcal A}^{(1)}_{\Lambda_k,-}(u)
\),
then
\[
\mathfrak X_{\Lambda_k}^{{\rm sel},-}
\big(u;\widehat{\mathcal A}^{(1)}_{\Lambda_k}\big)
=
\mathfrak X_{\Lambda_k}^{{\rm sel},-}
\big(u;\widetilde{\mathcal A}^{(1)}_{\Lambda_k}\big).
\]
\end{proposition}

\begin{proof}
Taking swap-odd parts of \eqref{eq:selector-same-scheme} gives
\[
\widehat{\mathcal A}^{(1)}_{\Lambda_k,-}(u)
-
\widetilde{\mathcal A}^{(1)}_{\Lambda_k,-}(u)
=
Q_tY^-_{\Lambda_k}(u).
\]
Let
\(
\Xi^-_{\Lambda_k}(u)
\in
\mathfrak X_{\Lambda_k}^{{\rm sel},-}
\big(u;\widetilde{\mathcal A}^{(1)}_{\Lambda_k}\big)
\).
Then
\[
Q_t\Xi^-_{\Lambda_k}(u)
=
\widetilde{\mathcal A}^{(1)}_{\Lambda_k,-}(u),
\]
so
\[
Q_t\big(
\Xi^-_{\Lambda_k}(u)+Y^-_{\Lambda_k}(u)
\big)
=
\widehat{\mathcal A}^{(1)}_{\Lambda_k,-}(u).
\]
Because both terms are \(\Gamma\)-odd, the sum lies in
\(
\mathfrak X_{\Lambda_k}^{{\rm sel},-}
\big(u;\widehat{\mathcal A}^{(1)}_{\Lambda_k}\big)
\).
This proves one inclusion, and the reverse inclusion follows by symmetry with \(-Y^-_{\Lambda_k}(u)\).

If \(Y'_{\Lambda_k}(u)\) is another primitive, then
\(
Q_t(Y'_{\Lambda_k}(u)-Y_{\Lambda_k}(u))=0
\).
Taking swap-odd parts yields
\[
Q_t\big(
Y_{\Lambda_k}^{\prime -}(u)-Y^-_{\Lambda_k}(u)
\big)=0.
\]
Since the difference is also \(\Gamma\)-odd, it lies in \(\mathcal K^-_{\Lambda_k}\).
If
\(
\widehat{\mathcal A}^{(1)}_{\Lambda_k,-}(u)
=
\widetilde{\mathcal A}^{(1)}_{\Lambda_k,-}(u)
\),
then one may choose \(Y^-_{\Lambda_k}(u)=0\), and the last claim follows.
\end{proof}

\begin{theorem}[Scheme-stable affine translation of the odd selector output]
\label{thm:selector-scheme-stability}
Let \(S\) and \(S'\) be cohomologically admissibly comparable finite-cutoff schemes on the same
chosen monotone sector.
Assume in addition that:
\begin{enumerate}[leftmargin=2em]
\item the canonical anomaly representatives
  \(
  \mathcal A^{(1)}_{k,S}(u)
  \)
  and
  \(
  \mathcal A^{(1)}_{k,S'}(u)
  \)
  are swap-even;
\item there exist chosen \(\Gamma\)-stable selector classes
  \[
  \mathcal C_{{\rm sel},\Lambda_k}^{S}(u),
  \qquad
  \mathcal C_{{\rm sel},\Lambda_k}^{S'}(u),
  \]
  and chosen selector representatives
  \[
  \widehat{\mathcal A}^{(1)}_{k,S}(u)\in\mathcal C_{{\rm sel},\Lambda_k}^{S}(u),
  \qquad
  \widehat{\mathcal A}^{(1)}_{k,S'}(u)\in\mathcal C_{{\rm sel},\Lambda_k}^{S'}(u),
  \]
  and an even sufficiently local cochain \(Y_k(u)\) such that
  \begin{equation}
  \widehat{\mathcal A}^{(1)}_{k,S'}(u)
  -
  \widehat{\mathcal A}^{(1)}_{k,S}(u)
  =
  \mathcal A^{(1)}_{k,S'}(u)
  -
  \mathcal A^{(1)}_{k,S}(u)
  +
  Q_tY_k(u).
  \label{eq:selector-exact-comparison}
  \end{equation}
\end{enumerate}
Define the odd comparison primitive
\[
Y^-_k(u)
:=
\frac12\Big(
Y_k(u)-\Gamma Y_k(u)
\Big).
\]
Then
\[
\mathfrak X_{{\Lambda_k},S'}^{{\rm sel},-}
\big(u;\widehat{\mathcal A}^{(1)}_{k,S'}\big)
=
\mathfrak X_{{\Lambda_k},S}^{{\rm sel},-}
\big(u;\widehat{\mathcal A}^{(1)}_{k,S}\big)
+
Y^-_k(u).
\]
If \(Y'_k(u)\) is another comparison primitive satisfying
\eqref{eq:selector-exact-comparison} for the same chosen pair of representatives, then
\[
Y_k^{\prime -}(u)-Y^-_k(u)\in \mathcal K^-_{\Lambda_k}.
\]
\end{theorem}

\begin{proof}
Taking swap-odd parts of \eqref{eq:selector-exact-comparison} and using the swap-evenness of the
canonical anomaly representatives gives
\[
\widehat{\mathcal A}^{(1)}_{k,S',-}(u)
- 
\widehat{\mathcal A}^{(1)}_{k,S,-}(u)
=
Q_tY^-_k(u).
\]
The same affine-fiber argument as in \cref{prop:selector-same-scheme} therefore gives the
translation statement by \(Y^-_k(u)\), and the same odd-kernel argument gives the ambiguity claim
for alternative primitives.
\end{proof}

\begin{remark}[Selector output handed to the downstream notes]
\label{rem:selector-handoff}
The raw selector datum is the affine class
\[
\mathfrak X_{\Lambda_k}^{{\rm sel},-}
\big(u;\widetilde{\mathcal A}^{(1)}_{\Lambda_k}\big)
\]
attached to the chosen selector representative.
Its ambiguity descent, transported dominant-response value, and weak-interface image are treated in
\cref{thm:dominant-quotient-descent,thm:active-response-transport-minimal,thm:char-weak-interface}.
\end{remark}

\begin{remark}[Relation to the Hodge-normalized layer]
\label{rem:selector-vs-hodge}
The Hodge-normalized comparison layer for admissible alternative representatives is treated in
\cref{uomdesc-thm:uomdesc-conditional-hodge-package}.
The present section uses only raw sufficiently local selector classes.
\end{remark}

\ifdefined\UOMTransportFragmentLoaded\else
\end{document}
\fi
